% Generated draft: requires mathematical and source review.
\documentclass[11pt]{article}
\usepackage[margin=25mm]{geometry}
\usepackage{fontspec}
\setmainfont{Latin Modern Roman}
\usepackage{amsmath,amssymb,mathtools,unicode-math}
\setmathfont{Latin Modern Math}
\usepackage{enumitem,needspace,xurl,hyperref}
\setlength{\parindent}{0pt}
\setlength{\parskip}{5pt}
\setlength{\emergencystretch}{3em}
\newcommand{\N}{\mathbb{N}}
\newcommand{\Z}{\mathbb{Z}}
\newcommand{\Q}{\mathbb{Q}}
\newcommand{\R}{\mathbb{R}}
\newcommand{\C}{\mathbb{C}}
\newcommand{\F}{\mathbb{F}}
\newcommand{\A}{\mathbb{A}}
\newcommand{\PP}{\mathbb{P}}
\newcommand{\E}{\mathbb{E}}
\newcommand{\eps}{\varepsilon}
\newcommand{\dd}{\,\mathrm{d}}
\begin{document}
\section*{UnsolvedMath ID: 30003579}
\textbf{Source problem code:} catalogue entry
\subsection{Definitions and mathematical statement}
The supplied record describes a research objective rather than a fully specified conjecture. Its stated ground fields are function fields. If global function fields are intended, such a field is a finite extension \(K/\F_q(t)\), where \(\F_q\) is a finite field. The record does not explicitly specify this global setting, a field \(K\), an algebraic group, a class of automorphic representations, or an associated \(L\)-function construction.

For clarity, equidistribution requires both a space of objects and a limiting probability measure. Given a topological space \(X\), probability measures \(\mu_T\) describing a family at parameter \(T\), and a proposed limiting probability measure \(\mu\), weak equidistribution means
\[
\lim_{T\to\infty}\int_X f(x)\,\dd\mu_T(x)
=\int_X f(x)\,\dd\mu(x)
\qquad\text{for every bounded continuous }f:X\to\R.
\]
This is a definition of one standard equidistribution formulation, not an additional hypothesis supplied by the source. Here the source does not determine \(T\), \(X\), \(\mu_T\), or \(\mu\). In particular, it supplies neither family weights nor a topology on a space of complete \(L\)-functions.

\emph{Source objective.} Establish equidistribution of local factors, or alternatively of complete \(L\)-functions, in families of automorphic representations over function fields subject to specified local conditions. Preferably, develop a method that also proves the Ramanujan conjecture for those families.

The local conditions, places being observed, normalization of local factors, associated \(L\)-functions, and family growth parameter remain unspecified. The word ``complete'' is retained without deciding whether it means completed \(L\)-functions or entire collections of Euler factors. The Ramanujan objective also requires a specified class of representations and a precise local assertion; a temperedness formulation would ask that the local components \(\pi_v\) be tempered at the places under consideration, but neither that scope nor the hypotheses needed for such a formulation are supplied. The preference for a common method is an additional research aim, not an asserted logical implication from equidistribution. Consequently, the record does not determine a unique quantified conjecture.
\subsection{Short English statement}
Seek equidistribution of local factors, or complete \(L\)-functions, in automorphic families over function fields with prescribed local conditions, preferably using a method that also proves the relevant Ramanujan conjecture. The families, limiting measures, normalizations, and precise Ramanujan assertion must be specified before this becomes a definite conjecture.
\subsection{Sources}
\begin{itemize}
\item Supplied local UnsolvedMath registry record: permanent ID \(30003579\), source problem code \texttt{catalogue entry}, title ``Equidistribution of Automorphic L Functions over Function Fields,'' in the Oberwolfach Reports Open Problems collection. No individual proposer is supplied; the proposed year is recorded as 2017.
\item Supplied report attribution: \emph{Automorphic Forms and Arithmetic (2018)}, Oberwolfach Reports, pp.\ 2475--2537, DOI \texttt{10.4171/owr/2017/40}. \url{https://doi.org/10.4171/owr/2017/40}. The collection describes Oberwolfach Reports as documenting mathematical workshops and their open-problem discussions. No problem-specific page or passage is supplied.
\item Supplied literature source: Blomer--Brumley, \url{https://doi.org/10.4007/annals.2011.174.1.18}. The record supplies no article title or problem-specific supporting passage.
\item Supplied literature source: OWR report, \url{https://doi.org/10.4171/owr/2017/40}.
\item Literature assessment dated exactly \texttt{2026-08-21}: the supplied curation labels the problem open and states that the Ramanujan conjecture is known in several cases, that strong bounds are available, and that the general automorphic statement over number fields remains open. This is supplied dated curation, not an independently verified assessment. Its number-field discussion does not specify the intended function-field problem. The background explicitly requests independent verification of time-sensitive status and cited claims.
\item The supplied statement status is \texttt{corrected\_verified}; its accompanying explanation says that mathematical content was preserved while notation, typography, and sentence structure were repaired for publication. This supplied editorial label does not establish mathematical correctness or independent source verification.
\end{itemize}
\end{document}
